% ---------------------------------------------------------------------------
% Section draft for The Web Conference (WWW) 2027 — ACM sigconf, double-blind.
% Companion to the MMSys 2027 paper "Enabling Browser Support of Arbitrary
% Multimedia via Dynamic WebAssembly-Based Pipelines". Numbers come from
% scripts/bench.js in the bevara-filter-manager repository (see the README
% in this directory to reproduce them).
%
% Before submission: check the anonymisation policy for the tool name
% (\toolname) and repository URLs, and re-run the benchmark on the
% evaluation machine so the figures match Table~\ref{tab:reg-publish}.
% ---------------------------------------------------------------------------
\newcommand{\toolname}{\texttt{bevara}}
\newcommand{\registry}{registry}

\section{Supplying Accessors through a Capability-Indexed, Verifiable Registry}
\label{sec:registry}

The pipeline mechanism of the previous sections deliberately knows nothing at
compile time about the accessors it will link: the set of decoders available
to a page is a property of the page (\texttt{with} attribute), of the site
(library policy of the \texttt{MutationObserver} and Service Worker modes), or
of the content itself (ISOBMFF encapsulation). All three supply policies
therefore assume that somebody \emph{curates} a collection of independently
compiled side modules, and the on-demand resolution of
Section~\ref{sec:ondemand} explicitly requires a mapping from a declared
stream capability to a side-module name. In practice this collection is
maintained by hand: a directory of \texttt{<name>\_<abi>.wasm} files copied
from a build tree, with no record of who produced them, from which sources, or
whether the copy in a page still matches what its author published. This
section fills that gap with a package model and a registry designed around
three properties that general-purpose WebAssembly registries do not offer: the
registry understands \emph{what} an accessor is (a GPAC filter with declared
capabilities, linked as a dynamic side module), it serves accessors directly to
browsers, and every publication is signed and appended to a log that consumers
can verify without trusting the registry.

\subsection{Package Model}
\label{sec:registry-model}

An accessor package consists of one WebAssembly side module and a manifest,
\texttt{bevara.json}, shown in Listing~\ref{lst:manifest}. Beyond the fields
that any package manager records (name, semantic version, description,
license, dependencies), the manifest carries the information that the
resolver of Section~\ref{sec:resolution} consumes: the GPAC filter names
registered by the module, a summary of its input and output capabilities
(stream type and codec), the \texttt{priority} used as a tie-breaker, and the
legal side conditions (\texttt{licence\_required}, \texttt{georestrictions})
by which a site may exclude accessors wholesale.

\begin{lstlisting}[caption={Manifest of the Theora accessor. Fields below the
blank line are build-time only and are stripped when publishing.},
label={lst:manifest},basicstyle=\ttfamily\scriptsize]
{ "name": "bevara:theora", "version": "1.2.0", "abi": 1,
  "description": "Theora decoder", "filters": ["theoradec"],
  "support": ["video"], "formats": ["Theora"],
  "inputs":  [{"stream_type":"visual","codec":"theora"}],
  "outputs": [{"stream_type":"visual","codec":"raw"}],
  "priority": 50, "license": "BSD-3-Clause",
  "licence_required": false, "georestrictions": false,
  "dependencies": {"bevara:ogg": "^1.0.0"},
  "sources": "https://github.com/.../theora",

  "main": "build/theora_1.wasm",
  "scripts": {"build": "emcmake cmake -S . -B build && cmake --build build"} }
\end{lstlisting}

\paragraph{Names and ABI.}
Package names carry a namespace, following the convention of the Bytecode
Alliance registry protocol~\cite{warg} (e.g.\ \texttt{bevara:theora}), so that packages remain
addressable if they are ever mirrored to a component registry. Each account
owns the namespace of its user name; the first publisher of a new namespace
claims it; a default namespace holds the reference accessors and is reserved
to administrators. The compiled artifact keeps the short file name
\texttt{<package>\_<abi>.wasm}, where the integer \texttt{abi} is the suffix that
the \texttt{with} attribute and the dynamic loader already use
(\texttt{theora\_1}), and it is bumped only when the host interface of the
main module changes, independently of the semantic version of the accessor
itself. Versions are immutable once published; owners may \emph{deprecate}
or \emph{yank} them, both of which are recorded (Section~\ref{sec:registry-trust}).

\paragraph{Validation at publication.}
The registry does not store opaque blobs. On upload it unpacks the tarball,
checks the manifest against the schema, and parses the binary: the module must
be well-formed and must carry a \texttt{dylink.0} custom
section~\cite{emscripten-dylink}, i.e.\ it must have been linked as a side
module that imports the shared memory and function table and exports
\texttt{\_\_wasm\_call\_ctors}. A module that merely \emph{compiles} to
WebAssembly but was linked as a standalone program is refused with an
explanation, since \texttt{dlopen()} in the main module could not link it.
The import and export tables extracted during this pass are stored with the
version and exposed by \texttt{\toolname{} info}, which gives site developers
a cheap check that, e.g., the expected \texttt{*\_register} entry points exist.

\subsection{Capability Index and Delivery to Browsers}
\label{sec:registry-index}

Two views of the catalogue are served for machines. A compact index,
\texttt{/accessors/index.json}, lists for every non-deprecated package its
latest version, file name, capabilities, priority, legal flags, size and
integrity digest; it is the ``mapping from a declared stream capability to a
side module name'' that deferred resolution needs
(Section~\ref{sec:ondemand}): when the demultiplexer of an Ogg container
declares a \texttt{visual/theora} PID with no accepting filter loaded, the
resolver selects, among index entries whose \texttt{inputs} match, the one with
the best priority and requests \texttt{theora\_1.wasm}. A second view
reproduces the \texttt{filter\_list.json} format of the existing filter
repository, so current builds of the wrapper work unchanged.

Accessors themselves are fetched from the registry like any static asset.
\texttt{/accessors/theora\_1.wasm} returns the newest non-deprecated version
for that ABI with a short cache lifetime, which is what a page using the
library policy wants. Under the same prefix, the versioned form
\texttt{bevara:theora/1.2.0/theora\_1.wasm} is instead immutable and cached
indefinitely, which is what a page with an explicit
\texttt{with} list and a lock file wants. Both are served with permissive
CORS and cross-origin resource policy headers so that any origin may link
them, and with a Brotli-compressed copy produced once at publication and
negotiated through \texttt{Accept-Encoding}: the 13 accessors of
Table~\ref{tab:reg-publish} shrink from 5.15\,MB to 1.36\,MB (26\%). Sites
that prefer to self-host run \texttt{\toolname{} install}, which resolves the
transitive closure of the declared dependencies (a decoder usually needs its
demultiplexer), downloads each module into the page's \texttt{accessors/}
directory under the exact name the \texttt{with} attribute expects, and prints
the resulting attribute value. Resolution is flat: because the loader
identifies modules by file name, a page can hold a single version of each
accessor, and conflicting ranges are reported rather than silently resolved.

\subsection{Verifiable Publication}
\label{sec:registry-trust}

An accessor executes with full access to the shared linear memory of the
pipeline, in every page that links it. The registry is thus a supply-chain
target in the sense documented for npm~\cite{ohm2020backstabber,
zimmermann2019smallworld}: an attacker who obtains a publisher's token, or who
compromises the registry host or its CDN, can substitute a module that will be
linked into thousands of pages. We address three threats: (T1) a stolen
credential used to publish under a maintainer's name; (T2) a registry that
serves different bytes for the same version to different clients, or replaces
a version after the fact; (T3) a registry that silently removes a version or a
package from its history. The design borrows the record-and-log structure of
Warg~\cite{warg} and of transparency logs in general~\cite{rfc6962,
sigstore2022}, but stays deliberately small so that it can be verified by a
few hundred lines of client code with no dependency beyond a hash function and
a signature scheme.

\paragraph{Signed records (T1).}
Each publisher holds an Ed25519~\cite{bernstein2012ed25519} key pair, created
by the CLI on first use and registered on the account. A publication is
described by a \emph{record}---a canonical JSON object with the package name,
version, ABI, file name, the SHA-256 of the module and of the tarball, the key
identifier and a timestamp---signed by the publisher's private key and sent
alongside the tarball. The registry recomputes the digests from the uploaded
bytes, checks that they equal those in the record, that the key is registered
to the authenticated account and not revoked, and that the signature verifies.
A token alone no longer suffices to publish, and a revoked key invalidates
future publications without affecting the validity of past ones.

\paragraph{Append-only log (T2).}
Accepted records are appended to a hash chain,
$h_i = \mathrm{SHA256}(h_{i-1} \,\|\, \mathrm{SHA256}(r_i) \,\|\, \sigma_i)$
with $h_0 = 0^{256}$, where $r_i$ is the canonical record and $\sigma_i$ its
signature. Yank and deprecation operations are recorded in the same chain,
signed by the registry's own key so that their origin is explicit. A
\emph{checkpoint} $\langle n, h_n, t\rangle$ signed by the registry key is
returned with every publication, every audit and with the capability index.
The chain is linear rather than a Merkle tree: inclusion of entry $i$ is
proven by the $n-i$ subsequent digests, which is $O(n)$ instead of
$O(\log n)$, but the log of a codec registry grows by a few entries per week,
verification is cheap (Table~\ref{tab:reg-log}), and the client needs no
tree logic. We consider this an appropriate trade-off for a self-hosted
registry; a deployment expecting millions of entries would switch to the
Merkle structure of Warg without changing the record format.

\paragraph{Client-side pinning (T3).}
Verification happens where the trust is needed: in the site. When
\texttt{\toolname{} install} downloads a module it checks the digest against
the version document, verifies the publisher signature over the record, and
writes to the lock file the digest, the key identifier and the log position of
every accessor. \texttt{\toolname{} audit} then fetches the log entries
published since the checkpoint pinned in the lock file, re-derives the chain
from the pinned head, checks the new checkpoint's signature with the pinned
registry key, and only then advances the pin. A registry that rewrote,
reordered or dropped an entry cannot produce a chain that extends a head the
client has already recorded, and a registry that presents different histories
to different clients is caught as soon as two of them compare pins---the same
gossip argument as for certificate transparency~\cite{rfc6962}. The audit also
re-hashes the entries it fetches for the installed versions rather than
trusting served hashes, flags versions that were yanked after installation,
and reports advisories (a semver range with a severity) that administrators
attach to packages, in the manner of \texttt{npm audit}.

\subsection{Tooling}
\label{sec:registry-tooling}

The CLI mirrors the vocabulary of npm so that web developers need not learn a
new workflow. For accessor authors, \texttt{init} scaffolds a filter project
(manifest, CMake rules that link a side module, a GPAC filter skeleton with
its capability table and self-registration constructor);
\texttt{build}, \texttt{pack} and \texttt{publish} compile, bundle, sign and
upload; \texttt{deprecate}, \texttt{unpublish} and \texttt{owner} maintain
published packages. For site developers, \texttt{install}, \texttt{update},
\texttt{outdated}, \texttt{uninstall} and \texttt{list} manage the
\texttt{accessors/} directory and the lock file, \texttt{audit} performs the
checks above, and \texttt{search}, \texttt{info} and \texttt{log} query the
catalogue and the transparency log. Authentication uses the registry's
identity provider through a browser round-trip that hands a token back to the
CLI on a loopback port, as \texttt{npm login} does. The registry is a single
process with an embedded database and a file store, packaged as a container
and deployed with a persistent volume; the whole system (registry, CLI, tests)
is about 3,000 lines of JavaScript with two third-party dependencies.

\subsection{Cost}
\label{sec:registry-eval}

We measured the registry and CLI on the laptop of Section~\ref{sec:setup}
(Apple M3 Pro), registry and client on the same machine, with the 13
accessors of the Showcase that have a manifest in the filter repository.
Table~\ref{tab:reg-publish} breaks down a publication. Packing and signing on
the client take a few milliseconds; on the server, parsing the module,
verifying the signature and appending to the log take 4.6\,ms in the median,
and the remainder is the one-off Brotli compression at quality 11, roughly
1.2\,ms per kilobyte of module. The capability index costs 611\,B per
accessor uncompressed (1.6\,kB Brotli for 13 accessors); extrapolated to the
86 accessors of Table~\ref{tab:eager}, the index a page must fetch to resolve
on demand is about 52\,kB uncompressed and 11\,kB compressed, against the
27.5\,MB that eager linking of the same 86 accessors transfers. Resolving a
three-accessor dependency graph takes under a millisecond of client time, and
downloading, hashing and signature-checking three modules totalling 1.6\,MB
takes 5\,ms over loopback, i.e.\ verification is negligible against transfer.

\begin{table}[t]
\caption{Publication cost per accessor (median of one run each; sizes in kB,
times in ms). \emph{Server} includes unpacking, manifest and side-module
validation, signature verification, storage and log append; \emph{brotli} is
the share of that time spent compressing the module at quality~11.}
\label{tab:reg-publish}
\centering\footnotesize
\setlength{\tabcolsep}{4pt}
\begin{tabular}{lrrrrrr}
\toprule
Accessor & wasm & brotli & pack & sign & server & \emph{brotli} \\
\midrule
theora        &   36 &  15 &  1.9 & 0.11 &   38 &   37 \\
ogg           &  214 &  71 &  5.2 & 0.24 &  269 &  266 \\
vorbis        &  127 &  50 &  3.8 & 0.16 &  173 &  173 \\
liba52        &  317 &  87 &  6.8 & 0.28 &  391 &  375 \\
libjxl        &  702 & 165 & 12.8 & 0.44 &  987 &  957 \\
ffmpeg-hevc   & 1236 & 261 & 22.5 & 0.67 & 1499 & 1484 \\
isobmff       & 1326 & 387 & 30.4 & 0.74 & 1608 & 1605 \\
\midrule
All 13 (total $|$ median) & 5151 & 1356 & 5.2 & 0.27 & 269 & 266 \\
\bottomrule
\end{tabular}
\end{table}

Table~\ref{tab:reg-log} gives the cost of verifying the log on the client,
for synthetic logs of $10^3$ and $10^4$ signed publications by 20 publishers.
Recomputing the hash chain alone costs 2.6\,\textmu s per entry; checking
every Ed25519 signature as well costs 78\,\textmu s per entry, so a full
verification from genesis of a ten-thousand-entry log---several years of
activity for a codec registry---takes 0.8\,s and transfers 6.5\,MB, while
the incremental audit a site performs routinely (50 new entries since its
pinned checkpoint) takes 4\,ms and 33\,kB. Verification therefore fits in a
continuous-integration step or a pre-deployment hook without noticeable cost.

\begin{table}[t]
\caption{Client-side verification of the transparency log (single core).}
\label{tab:reg-log}
\centering\small
\begin{tabular}{lrr}
\toprule
 & $N = 10^3$ & $N = 10^4$ \\
\midrule
Log size (JSON)                 & 648\,kB & 6.5\,MB \\
Hash chain only                 & 2.9\,ms & 26\,ms \\
Chain + all signatures          & 79\,ms  & 776\,ms \\
Incremental (50 new entries)    & 4.2\,ms & 4.0\,ms \\
\bottomrule
\end{tabular}
\end{table}

\subsection{Relation to Existing Registries}
\label{sec:registry-related}

npm~\cite{npm} established the workflow we mirror (manifest, semantic
versioning, lock file, \texttt{audit}) and its ecosystem studies motivate the
threat model~\cite{zimmermann2019smallworld,ohm2020backstabber}, but npm
packages are JavaScript and its registry neither validates WebAssembly nor
serves it to browsers. WAPM~\cite{wapm}, now the Wasmer registry, distributes
WASI programs for a server-side runtime; the Bytecode Alliance registry
protocol Warg~\cite{warg} and the \texttt{wkg} tooling distribute
\emph{components} with typed WIT interfaces, increasingly over OCI
registries~\cite{oci-artifacts}. Neither models the object we distribute: a
side module with a C ABI towards a host that shares linear memory with it, a
choice the component model excludes by design and on which the zero-copy
pipelines of Section~\ref{sec:performance} depend. What we retain from Warg
is its naming and its record-and-log structure, and from
Sigstore~\cite{sigstore2022} and TUF~\cite{samuel2010tuf} the principle that
consumers, not the registry, are the verifiers. Should the component model
gain a dynamic-linking story with shared memory, accessors could become
components and Warg their transport; the capability index and the browser
delivery layer described here would remain the part specific to media
pipelines.
